\subsection{微分した商を調べてよい条件}

極限を計算していると，分子も分母も$0$へ近づく商が現れる。その場合，分子と分母がどれだけ速く$0$へ近づくかを比べる必要がある。微分係数がその速さを表すことから，分子と分母の微分を比べる方法が考えられる。ただし，元の商の形を確かめず，機械的に分子と分母を微分してはいけない。

\begin{theorem}{ロピタルの定理の一側極限版}
$x$が区間の一端$c$へ近づく一側極限を考える。$c$は有限でも$+\infty,-\infty$でもよい。端点を除くその区間で$f,g$は微分可能であり，$g'(x)\ne0$とする。さらに，次のいずれかを満たすとする。
\begin{enumerate}[label=(\roman*),leftmargin=*]
 \item $f(x)\to0$かつ$g(x)\to0$である。
 \item $|f(x)|\to\infty$かつ$|g(x)|\to\infty$である。
\end{enumerate}
このとき
\[
 \lim_{x\to c}\frac{f'(x)}{g'(x)}=L
\]
が存在すれば
\[
 \lim_{x\to c}\frac{f(x)}{g(x)}=L
\]
である。ここで$L$は実数のほか，$+\infty,-\infty$でもよい。元の商は端点の十分近くで定義されている。
\end{theorem}

両側極限を求めるときは，左右の一側極限にそれぞれ適用し，結果が一致することを確認する。また，$\frac{f'}{g'}$をさらに微分して計算したい場合は，その新しい商についても同じ条件を改めて確認する。微分後の商が$0$と$0$の商になるとは限らない。

